% SPE Final Project — Drift-Aware Churn Prediction Pipeline
% Compile with: pdflatex (run twice for TOC). Optional: latexmk -pdf spe_final_report.tex
\documentclass[11pt,a4paper,oneside]{report}

\usepackage[utf8]{inputenc}
\usepackage[T1]{fontenc}
\usepackage{lmodern}
\usepackage[english]{babel}
\usepackage{geometry}
\geometry{margin=1in}
\usepackage{graphicx}
\usepackage{booktabs}
\usepackage{hyperref}
\usepackage{url}
\usepackage{caption}
\usepackage{enumitem}
\usepackage{parskip}
\usepackage{microtype}

\hypersetup{
  colorlinks=true,
  linkcolor=blue!60!black,
  citecolor=blue!60!black,
  urlcolor=blue!60!black,
  pdftitle={SPE Final Project: Drift-Aware Churn Prediction Pipeline},
}

% --- Title page: edit the macros below or replace the titlepage block entirely ---
\newcommand{\ReportTitle}{Drift-Aware Churn Prediction Pipeline}
\newcommand{\ReportSubtitle}{Software Platform Engineering Final Project}
\newcommand{\UniversityName}{[University / Institution Name]}
\newcommand{\DepartmentName}{[Department Name]}
\newcommand{\AuthorName}{[Your Name]}
\newcommand{\StudentId}{[Student ID]}
\newcommand{\CourseCode}{[Course Code]}
\newcommand{\ReportDate}{\today}

% Placeholder image: framed box (replace with \includegraphics when you add screenshots)
\newcommand{\screenshotplaceholder}[2]{%
  \begin{center}
    \fbox{%
      \parbox[c][0.22\textheight][c]{0.92\linewidth}{%
        \centering\small\textit{[Screenshot placeholder: #1]}\\[0.5em]
        #2%
      }%
    }%
  \end{center}%
}

\begin{document}

% ==================== TITLE PAGE ====================
\begin{titlepage}
  \centering
  \vspace*{1.5cm}
  {\Large\bfseries \UniversityName\par}
  \vspace{0.5cm}
  {\large \DepartmentName\par}
  \vspace{2cm}
  {\LARGE\bfseries \ReportTitle\par}
  \vspace{0.75cm}
  {\large \ReportSubtitle\par}
  \vspace{0.5cm}
  {\large \CourseCode\par}
  \vfill
  {\large \textbf{\AuthorName}\par}
  \vspace{0.25cm}
  {\normalsize \StudentId\par}
  \vspace{1.5cm}
  {\large \ReportDate\par}
  \vspace*{1cm}
\end{titlepage}

\pagenumbering{roman}
\setcounter{page}{2}

\begin{abstract}
\noindent
This report describes a software platform engineering project implementing an end-to-end \textbf{drift-aware customer churn prediction} system for a banking-style use case. Four \textbf{Flask} microservices handle data ingestion, model training, online prediction, and statistical drift detection against a reference distribution saved at training time. The platform runs in \textbf{Docker} and \textbf{Kubernetes}, with a shared persistent volume for model artifacts, optional \textbf{horizontal pod autoscaling}, \textbf{Jenkins} continuous integration and deployment, optional \textbf{Ansible} provisioning with \textbf{Vault}-encrypted secrets, and an \textbf{ELK} stack (\textbf{Elasticsearch}, \textbf{Filebeat}, \textbf{Kibana}) for structured logs and dashboards. When drift is detected, the drift service can optionally trigger retraining by calling the training API. The report summarizes requirements, machine learning design, service APIs, deployment topology, CI/CD, observability, security considerations, and evaluation scenarios. Screenshots are indicated by placeholders throughout; replace them with your captures where noted.
\end{abstract}

\tableofcontents
\clearpage
\pagenumbering{arabic}

% ==================== CHAPTER 1 ====================
\chapter{Introduction}

\section{Background and motivation}
Banks and financial institutions rely on machine learning to estimate \textbf{customer churn}: the risk that an account holder will leave. Models trained on historical data can become unreliable when the live data distribution shifts (for example, changing demographics, products, or economic conditions). Production ML systems therefore need clear \textbf{operational practices}: modular services, automated testing and deployment, observability, and explicit handling of \textbf{data drift}.

\section{Problem statement}
The project addresses churn prediction from tabular customer records while ensuring that \textbf{deployment and monitoring} are first-class concerns. The system must accept batches of customer rows, score churn, and compare incoming batches to a \textbf{reference distribution} established during training so that drift can be flagged and, optionally, \textbf{model retraining} can be triggered automatically.

\section{Project goals}
\begin{itemize}[leftmargin=*]
  \item Deliver four cooperating \textbf{microservices} with stable HTTP APIs and automated tests.
  \item Package and orchestrate services with \textbf{Docker} and \textbf{Kubernetes}, including shared storage for artifacts.
  \item Implement \textbf{CI/CD} with \textbf{Jenkins} (build, test, deploy, and ELK-related steps as defined in the pipeline).
  \item Support optional \textbf{Ansible}-based provisioning and \textbf{ELK}-based log aggregation and visualization.
  \item Emit \textbf{structured JSON logs} suitable for search and dashboards in Kibana.
\end{itemize}

\section{Scope and assumptions}
The implementation focuses on a single churn use case with a fixed feature schema, a \textbf{scikit-learn} logistic regression pipeline, and drift tests based on saved reference statistics. Multi-tenant isolation, advanced authn/z for APIs, and full production SLO/SLA guarantees are out of scope unless extended separately.

\begin{figure}[htbp]
  \screenshotplaceholder{High-level system or repository overview}{Optional: landing view of the repository README, project board, or a one-slide overview you use for demos.}
  \caption{Visual overview of the project context (optional). Replace the placeholder with your screenshot; useful for readers who have not seen the repository.}
  \label{fig:project-overview}
\end{figure}

\section{Report structure}
Chapter~\ref{ch:requirements} states requirements and a system overview. Chapter~\ref{ch:ml} covers ML and drift design. Chapter~\ref{ch:microservices} details microservices and APIs. Chapter~\ref{ch:platform} describes containers, Kubernetes, and Ansible. Chapter~\ref{ch:cicd} covers CI/CD and testing. Chapter~\ref{ch:observability} addresses observability and operations. Chapter~\ref{ch:security} summarizes security and governance. Chapter~\ref{ch:results} presents results and evaluation. Chapter~\ref{ch:conclusion} concludes and suggests future work.

% ==================== CHAPTER 2 ====================
\chapter{Requirements and system overview}
\label{ch:requirements}

\section{Functional requirements}
\begin{itemize}[leftmargin=*]
  \item \textbf{Ingestion:} Accept validated JSON batches matching the churn schema and orchestrate calls to serving and drift services.
  \item \textbf{Training:} Fit a preprocessing and classification pipeline, write \texttt{churn\_model.pkl} and \texttt{reference\_distribution.pkl} to shared storage, and expose health for orchestration.
  \item \textbf{Serving:} Load the trained model and return predictions and churn probabilities for rows that match the expected feature set.
  \item \textbf{Drift detection:} Compare each batch to the reference; return \texttt{drift\_detected} and details; optionally \texttt{POST} to the training URL when drift is detected and configured.
\end{itemize}

\section{Non-functional requirements}
Services should support \textbf{rolling updates} and readiness probes where configured, \textbf{horizontal scaling} for selected workloads (per deployment policy), structured logging for troubleshooting, and safe handling of \textbf{secrets} via Ansible Vault for automation.

\section{High-level architecture}
Ingestion invokes serving and drift. Training, serving, and drift read or write model-related files on a \textbf{shared PVC}. Drift may call training when drift is detected and \texttt{TRAINING\_URL} is set. Log shipping sends container JSON logs to Elasticsearch for Kibana.

\begin{figure}[htbp]
  \screenshotplaceholder{Architecture diagram}{Paste a screenshot of your architecture diagram (for example Mermaid export, draw.io, or a slide). Should show ingestion, training, serving, drift, PVC, and ELK path.}
  \caption{End-to-end architecture of the churn pipeline and observability path. Replace with your diagram screenshot.}
  \label{fig:architecture}
\end{figure}

\section{Technology stack}
Table~\ref{tab:stack} summarizes the main technologies.

\begin{table}[htbp]
  \centering
  \caption{Technology stack summary.}
  \label{tab:stack}
  \begin{tabular}{@{}ll@{}}
    \toprule
    \textbf{Layer} & \textbf{Technology} \\
    \midrule
    APIs & Python Flask \\
    ML & pandas, NumPy, SciPy, scikit-learn \\
    Containers & Docker \\
    Orchestration & Kubernetes (Deployments, Services, PV/PVC, HPA) \\
    CI/CD & Jenkins (\texttt{Jenkinsfile}) \\
    IaC / provisioning & Ansible (optional), Vault for secrets \\
    Logs / dashboards & Elasticsearch, Filebeat, Kibana \\
    \bottomrule
  \end{tabular}
\end{table}

% ==================== CHAPTER 3 ====================
\chapter{Machine learning design}
\label{ch:ml}

\section{Dataset and features}
The churn use case expects tabular rows with identifiers and attributes such as credit score, geography, gender, age, tenure, balance, product count, card and activity flags, estimated salary, and label \texttt{Exited} where applicable. Ingestion validates the \textbf{full} schema required by the project.

\section{Model}
Training builds a \textbf{scikit-learn} pipeline: preprocessing for categorical and numeric features followed by \textbf{logistic regression}. The fitted pipeline is persisted as \texttt{churn\_model.pkl} on the shared volume.

\section{Reference distribution}
At training time, summary information needed for drift checks is stored in \texttt{reference\_distribution.pkl}. The drift service loads this file to compare new batches against the training-time reference.

\section{Drift detection method}
Drift detection compares incoming batches to the reference using per-feature statistical tests (KS-style comparisons against synthetic draws from reference mean and standard deviation), plus label distribution drift on \texttt{Exited}. The API returns a boolean \texttt{drift\_detected} and a \texttt{details} object for inspection.

\section{Retraining policy}
When \texttt{drift\_detected} is true and the drift service has \texttt{TRAINING\_URL} configured, it can \textbf{POST} the batch JSON to the training service's \texttt{/train} endpoint so the model and reference can be refreshed. Ingestion always invokes drift as part of \texttt{/ingest}, so retraining is tied to the ingestion path when drift fires and retraining is enabled.

\begin{figure}[htbp]
  \screenshotplaceholder{Drift or training output}{Example: terminal or JSON showing \texttt{drift\_detected}, \texttt{details}, or \texttt{training\_run\_*} log lines.}
  \caption{Example evidence of drift detection or training lifecycle (JSON response or structured log). Replace with your screenshot.}
  \label{fig:drift-example}
\end{figure}

% ==================== CHAPTER 4 ====================
\chapter{Microservices and APIs}
\label{ch:microservices}

\section{Data ingestion service}
\texttt{POST /ingest} accepts a JSON array or single object with the full churn schema. It forwards to the configured serving and drift URLs and returns a combined response including \texttt{drift\_response} (mirroring drift JSON, including \texttt{drift\_detected}).

\section{Model training service}
\texttt{POST /train} trains on JSON rows with required columns (including \texttt{Exited}, \texttt{Geography}, \texttt{Gender}) or falls back to \texttt{train.csv} on the PVC. \texttt{GET /health} supports compose and Kubernetes health checks. Structured events include \texttt{training\_run\_started}, \texttt{training\_run\_completed}, and \texttt{training\_run\_failed}.

\section{Model serving service}
\texttt{POST /predict} accepts one row or an array; drops \texttt{Exited} if present; returns predictions using the model loaded from the PVC.

\section{Drift detection service}
\texttt{POST /drift} compares the batch to \texttt{reference\_distribution.pkl}. If the reference is missing, the API returns a clear JSON error. On success, the response includes \texttt{drift\_detected} and \texttt{details}; when a retrain call is made, \texttt{training\_response} may appear.

\section{Inter-service communication}
Downstream URLs are provided by environment variables (for example \texttt{SERVING\_URL}, \texttt{DRIFT\_URL}, \texttt{TRAINING\_URL}). Ports default to 5000 (ingest), 5001 (train), 5002 (serve), and 5003 (drift) in local documentation.

\begin{figure}[htbp]
  \screenshotplaceholder{API test}{Example: \texttt{curl}, Postman, or HTTP client showing a successful \texttt{/ingest} or \texttt{/predict} response.}
  \caption{Sample API interaction demonstrating a successful request and response body. Replace with your screenshot.}
  \label{fig:api-sample}
\end{figure}

% ==================== CHAPTER 5 ====================
\chapter{Platform and deployment}
\label{ch:platform}

\section{Containerization}
Each microservice is built as a Docker image with its own Dockerfile and dependencies. Images are tagged and pushed as part of the CI pipeline for deployment to the cluster.

\section{Kubernetes resources}
Manifests define Deployments and ClusterIP Services for each workload. A PersistentVolume and PersistentVolumeClaim mount a shared path (for example \texttt{/data/churn-model}) into training, serving, and drift pods. Horizontal Pod Autoscaler manifests apply to selected services (training is typically excluded from HPA where appropriate).

\begin{figure}[htbp]
  \screenshotplaceholder{Kubernetes resources}{Example: \texttt{kubectl get pods,svc,pvc,hpa} output or Kubernetes Dashboard view.}
  \caption{Kubernetes cluster state showing pods, services, PVC, and optionally HPA. Replace with your screenshot.}
  \label{fig:k8s-state}
\end{figure}

\section{Networking and service discovery}
Services communicate via in-cluster DNS names and ClusterIP ports. Clients access the stack through port-forward, ingress (if configured), or load balancers depending on your environment.

\section{Configuration and secrets}
Runtime configuration uses environment variables. Ansible variables for sensitive data can be encrypted with \textbf{Ansible Vault}.

\section{Infrastructure as code}
Ansible roles (for example Kubernetes and ELK-related tasks) support repeatable provisioning. Playbooks such as \texttt{site.yml} orchestrate role application.

\begin{figure}[htbp]
  \screenshotplaceholder{Ansible or inventory}{Example: Ansible playbook run output, directory layout, or Vault workflow (redacted).}
  \caption{Ansible automation evidence (run excerpt or role structure). Replace with your screenshot; redact secrets.}
  \label{fig:ansible}
\end{figure}

% ==================== CHAPTER 6 ====================
\chapter{CI/CD and quality assurance}
\label{ch:cicd}

\section{Jenkins pipeline}
The \texttt{Jenkinsfile} defines stages: for example Ansible checks, unit tests, Docker Compose integration tests, image build and registry push, Kubernetes deployment, and ELK setup. Exact stage names and order should match your repository.

\begin{figure}[htbp]
  \screenshotplaceholder{Jenkins pipeline}{Jenkins Blue Ocean or classic UI showing stages turning green for a successful build.}
  \caption{Jenkins pipeline run showing stages (Ansible, tests, compose, build, deploy, ELK). Replace with your screenshot.}
  \label{fig:jenkins}
\end{figure}

\section{Automated testing}
Each service includes unit tests under its package. Repository-level integration tests (for example \texttt{tests/test\_integration.py}) exercise the pipeline against a compose stack (\texttt{docker-compose.test.yml}) where applicable.

\begin{figure}[htbp]
  \screenshotplaceholder{Test results}{Example: pytest or Jenkins test report showing passed tests.}
  \caption{Automated test results (local pytest or Jenkins test stage). Replace with your screenshot.}
  \label{fig:tests}
\end{figure}

\section{Release and deployment flow}
A typical flow is: commit to version control, Jenkins builds and tests artifacts, images are pushed to a registry, manifests apply workloads to Kubernetes, and optional jobs configure Elasticsearch index patterns or Kibana dashboards.

% ==================== CHAPTER 7 ====================
\chapter{Observability and operations}
\label{ch:observability}

\section{Structured logging}
Each service emits JSON logs with fields suitable for Filebeat shipping to Elasticsearch, including service identifiers and training lifecycle events for traceability.

\section{ELK stack}
Filebeat collects container logs into Elasticsearch index patterns such as \texttt{project-logs-*}. Kibana provides search and dashboards for operational visibility.

\begin{figure}[htbp]
  \screenshotplaceholder{Kibana Discover}{Kibana Discover filtered to your service or \texttt{training\_run\_*} messages.}
  \caption{Kibana Discover showing structured logs from one or more microservices. Replace with your screenshot.}
  \label{fig:kibana-discover}
\end{figure}

\begin{figure}[htbp]
  \screenshotplaceholder{Kibana dashboard}{Optional: imported dashboard for log volume or errors by service.}
  \caption{Kibana dashboard (optional) for log or error visualization. Replace with your screenshot.}
  \label{fig:kibana-dashboard}
\end{figure}

\section{Monitoring and troubleshooting}
Operators verify pod readiness, inspect logs in Kibana for failed ingest or training, and confirm drift and serving URLs when debugging cross-service calls.

\section{Operational scenarios}
A heavy drift batch should yield \texttt{drift\_detected: true}; with \texttt{TRAINING\_URL} set, drift triggers retraining and the PVC receives updated artifacts. Logs should show the corresponding training events and downstream ingest behavior.

% ==================== CHAPTER 8 ====================
\chapter{Security and governance}
\label{ch:security}

\section{Secrets management}
Ansible Vault encrypts sensitive variables used in automation. Decryption keys are not stored in the repository.

\section{Network and access}
Internal service URLs should not be exposed beyond the cluster without an appropriate ingress and authentication strategy.

\section{Data and model lifecycle}
Training can be triggered by automation or drift callbacks; access to the training endpoint and PVC should be restricted in production environments.

% ==================== CHAPTER 9 ====================
\chapter{Results and evaluation}
\label{ch:results}

\section{Functional validation}
The README describes scenarios: a batch similar to training data should often yield \texttt{drift\_detected: false}, while an extreme batch should tend toward \texttt{drift\_detected: true}. Repeatability note: random sampling in drift checks can cause rare variability; larger batches reduce noise.

\section{End-to-end behavior}
Through \texttt{/ingest}, \texttt{drift\_response} mirrors \texttt{/drift}. When retraining is configured, heavy drift may include \texttt{training\_response} nested under drift output.

\section{Limitations}
The drift methodology uses reference summaries and statistical tests suited to the project scope; it is not a full production MLOps platform. Reference randomness may cause occasional borderline results on small batches.

\begin{figure}[htbp]
  \screenshotplaceholder{Before/after drift}{Side-by-side or two captures: \texttt{drift\_detected: false} vs \texttt{true} from curl or browser.}
  \caption{Comparison of no-drift versus heavy-drift responses (terminal or saved JSON). Replace with your screenshot.}
  \label{fig:drift-compare}
\end{figure}

% ==================== CHAPTER 10 ====================
\chapter{Conclusion and future work}
\label{ch:conclusion}

\section{Summary}
The project delivers a drift-aware churn prediction platform with clear separation of concerns across four Flask services, Kubernetes deployment, Jenkins CI/CD, optional Ansible provisioning, and ELK-based observability.

\section{Lessons learned}
Platform engineering adds operational complexity (manifests, pipelines, secrets) but improves repeatability and visibility compared to ad hoc scripts.

\section{Future improvements}
Possible extensions include stronger authentication for APIs, Prometheus metrics, canary deployments, experiment tracking, feature stores, and richer drift monitoring (for example population stability index or model-specific metrics).

\appendix
\chapter{Build and compile notes}
To compile this report: \texttt{pdflatex spe\_final\_report.tex} twice, or use \texttt{latexmk -pdf spe\_final\_report.tex}. To insert real figures, replace \verb|\screenshotplaceholder{...}{...}| blocks with standard \verb|\includegraphics| calls and keep the same \verb|\caption| labels for cross-references.

\end{document}
